%%
%% Speech Analytics and Natural Language Processing Applied to
%% Clinical Communication: A Scoping Review
%%
%% [Author identifiers removed for single-anonymized peer review]
%%
%% Target journal: Artificial Intelligence in Medicine (Elsevier)
%% Article type: Methodological Review
%% Class: elsarticle.cls (numbered/Vancouver-style references)
%%

\documentclass[review,12pt]{elsarticle}

\usepackage[numbers,sort&compress]{natbib}
\usepackage[utf8]{inputenc}
\usepackage[T1]{fontenc}
\usepackage{lmodern}
\usepackage[english]{babel}
\usepackage{amsmath}
\usepackage{amssymb}
\DeclareSymbolFont{CMSYfix}{OMS}{cmsy}{m}{n}
\DeclareMathSymbol{\approx}{\mathrel}{CMSYfix}{"19}
\usepackage{csquotes}
\usepackage{graphicx}
\usepackage{booktabs}
\usepackage{longtable}
\usepackage{array}
\usepackage{hyperref}
\usepackage{xcolor}
\usepackage{tikz}

\hypersetup{colorlinks=true, linkcolor=blue, citecolor=blue, urlcolor=blue}

\journal{Artificial Intelligence in Medicine}

%% ================================================================
\begin{document}

\begin{frontmatter}

%% ----------------------------------------------------------------
%% Title
%% ----------------------------------------------------------------
\title{Speech Analytics and Natural Language Processing Applied to
Clinical Communication: A Scoping Review}

%% ----------------------------------------------------------------
%% Authors — identifiers removed for single-anonymized review
%% ----------------------------------------------------------------
\author[inst1]{Author 1\corref{cor1}\fnref{fn1}}
\author[inst1]{Author 2\fnref{fn2}}
\author[inst2]{Author 3\fnref{fn3}}
\author[inst2]{Author 4\fnref{fn4}}
\author[inst2]{Author 5\fnref{fn5}}

\cortext[cor1]{Corresponding author.}
\fntext[fn1]{[Role and affiliation removed for review]}
\fntext[fn2]{[Role and affiliation removed for review]}
\fntext[fn3]{[Role and affiliation removed for review]}
\fntext[fn4]{[Role and affiliation removed for review]}
\fntext[fn5]{[Role and affiliation removed for review]}

\address[inst1]{[Institution 1 removed for review], Brazil}
\address[inst2]{[Institution 2 removed for review], Brazil}

%% ----------------------------------------------------------------
%% Abstract
%% ----------------------------------------------------------------
\begin{abstract}
\textbf{Background:}
Professional-patient communication constitutes a dense source of clinical
information that remains predominantly unstructured and analytically
underutilised. Recent advances in Automatic Speech Recognition (ASR),
Natural Language Processing (NLP) and vocal biomarker analysis have expanded
the capacity to extract structured data from clinical interactions during
anamnesis. Nevertheless, few studies operationalise these interactions as
measurable variables derived from speech, particularly with regard to
communicational and behavioural markers that can be used for data analysis
and outcome assessment in patient care.

\textbf{Objective:}
To identify, synthesise and critically evaluate the evidence on the application
of \textit{Speech Analytics} and NLP to professional-patient clinical
communication, mapping methodological approaches, identified markers, technical
challenges and gaps for future research.

\textbf{Methods:}
Scoping review following the PRISMA-ScR guidelines \citep{Tricco2018}
and JBI \citep{Peters2020}. Searches were conducted across five electronic
databases (PubMed/MEDLINE, Google Scholar, EBSCO/CINAHL, Web of Science, and
Cochrane Library; January 2014--July 2026; English, Spanish, and Portuguese),
using six search strings in PubMed/MEDLINE and five in the remaining databases,
and supplemented by manual \textit{snowballing}. Of 5,621 records identified
across the five databases (3,189 after deduplication) and 7 additional
records identified by \textit{snowballing}, 24~articles met the eligibility
criteria and were organised into four thematic axes: (A)~Clinical NLP
and biomedical entity extraction; (B)~automated clinical documentation and ASR;
(C)~physician-patient communication; (D)~vocal biomarkers and multimodal
analysis.

\textbf{Results:}
Studies evidence consolidated advances in clinical NLP (\textit{transformer}-based
models, BERT leads Named Entity Recognition) and automated transcription
(WER $\approx$9\% achievable with Whisper/PyAnnote). Patient-centred
communication demonstrates measurable impact on objective health outcomes.
Communicational markers -- self-repair, engagement, emotional support and
barrier indicators -- are identified as adherence predictors. Multimodal
analysis (acoustic\,+\,semantic) remains in an early stage; most approaches
focus on documentation automation, neglecting clinical interaction as an
object of analysis in its own right.

\textbf{Conclusions:}
A consistent gap is confirmed: few studies operationalise clinical interactions
as measurable variables derived from speech, particularly communicational
and behavioural markers extracted from the anamnesis. This review grounds the
need for studies integrating \textit{Speech Analytics} with structured analysis
of clinical communication as a data source for adherence monitoring and quality
of care.

\end{abstract}

%% ----------------------------------------------------------------
%% Keywords
%% ----------------------------------------------------------------
\begin{keyword}
Speech analytics \sep
Natural language processing \sep
Clinical communication \sep
Automatic speech recognition \sep
Communication markers \sep
Digital scribe
\end{keyword}

\end{frontmatter}

%% ================================================================
\section{Introduction}
\label{sec:intro}

Clinical communication is one of the richest and least explored data sources
in healthcare. In public health systems, anamnesis sessions and therapeutic
guidance -- verbal interactions with a propedeutic structure defined
by clinical semiology \citep{Porto2019} -- occur predominantly
as analytically unstructured exchanges, rarely converted into measurable
data \citep{Wang2018,Koleck2019}.

The quality of the professional-patient interaction has a measurable impact
on adherence and objective clinical outcomes \citep{Kelley2014}. Communication
barriers -- including health literacy gaps, misaligned expectations and
miscommunication patterns -- are associated with low therapeutic adherence
and worse prognosis \citep{McCabe2018,Alkhamees2023}.

\textit{Speech Analytics} -- the computational analysis of spoken language --
offers a pathway to extract structured and measurable indicators from these
interactions. Recent advances in ASR, particularly transformer-based models
such as Whisper \citep{Radford2023}, combined with instruction-tuned large
language models (LLMs), have lowered the barriers for deploying end-to-end
NLP \textit{pipelines} in healthcare \citep{Klusty2025}. Nevertheless, most
existing approaches target diagnostic support or documentation automation,
without treating the clinical interaction itself as an object of structured
analysis \citep{vanBuchem2021}. A prospective observational protocol for
monitoring orthopedic rehabilitation in SUS through \textit{Speech
Analytics} has recently been proposed [Author et al., 2026]. That protocol
acknowledges, as a central methodological requirement, the absence of a
consolidated evidence map on the operationalisation of clinical interaction
as a source of speech-derived measurable variables. The present scoping
review responds directly to that requirement, providing the systematic
evidence synthesis necessary to anchor the design of the prospective
empirical studies constituting Phase 2 of that research programme.

Recent reviews have addressed aspects partially overlapping with this scope.
\citet{Falcetta2023} identified the absence of prospective validation in
automated documentation systems for clinical interactions. \citet{Ng2025}
evaluated the technical performance of ASR systems for clinical documentation,
reporting WER values between 8.7\% and over 50\% depending on the context.
\citet{Alabbad2026} mapped ASR applications in healthcare in the post-LLM
era. \citet{Albert2024} conducted a scoping review on AI, speech and NLP
methods for clinician-patient communication assessment. None of these reviews,
however, focuses on the operationalisation of the clinical interaction as a
set of \textbf{speech-derived measurable variables} -- particularly
communicational and behavioural markers (self-repair, engagement, barriers,
health literacy) -- as an object of analysis in its own right, independent of
documentation or diagnosis. It is precisely this gap that the present review
aims to map.

This scoping review aims to map the evidence on \textit{Speech Analytics}
and NLP applied to clinical communication, characterising the methodological
landscape and identifying gaps for future research.

%% ================================================================
\section{Methods}
\label{sec:metodos}

\subsection{Study design}

This scoping review follows the methodological framework proposed by
\citet{Arksey2005}, refined by \citet{Levac2010} and the Joanna Briggs
Institute (JBI) \citep{Peters2020}. Conduct and reporting follow the
PRISMA-ScR guidelines (\textit{Preferred Reporting Items for Systematic
Reviews and Meta-Analyses} extension for Scoping Reviews) \citep{Tricco2018}.

\subsection{Research questions}

\begin{itemize}
  \item \textbf{Q1:} Which \textit{Speech Analytics} approaches have been
        applied to clinical communication over the past five years?
  \item \textbf{Q2:} Which NLP and ASR methods are used for automated
        transcription and analysis of clinical consultations?
  \item \textbf{Q3:} Which communicational and behavioural markers have been
        identified or proposed in the professional-patient interaction?
  \item \textbf{Q4:} What are the technical, ethical and methodological
        challenges for automated analysis of clinical communication?
\end{itemize}

\subsection{Eligibility criteria}

\subsubsection{Inclusion criteria}
\begin{itemize}
  \item \textbf{IC1:} The study applies \textit{Speech Analytics}, NLP or ASR
        technologies for clinical transcription and extraction of unstructured
        medical data, or analyses and measures the foundations of
        professional-patient communication.
  \item \textbf{IC2:} The study focuses on free-text EHR interactions, spoken
        clinical dialogues (consultations), automated documentation generation
        (digital scribes/SOAP notes) or theoretical frameworks of shared
        decision-making.
  \item \textbf{IC3:} Full-text availability in English, Portuguese or Spanish;
        published between January 2014 and 2026.
  \item \textbf{IC4:} Presents original development, empirical feasibility
        testing, algorithm validation, or constitutes a prior systematic or
        scoping review on the topic.
\end{itemize}

\subsubsection{Exclusion criteria}
\begin{itemize}
  \item \textbf{EC1:} Studies that do not address NLP, ASR or audio analysis
        of clinical interactions.
  \item \textbf{EC2:} Studies describing transcription or extraction systems
        without analysing interaction quality or impact on the
        professional-patient relationship.
  \item \textbf{EC3:} Exclusive focus on voice-based diagnosis of physical or
        psychiatric conditions without analysis of dialogic interaction.
  \item \textbf{EC4:} Opinion pieces, commentaries, editorials, conference
        abstracts or books.
  \item \textbf{EC5:} Studies that do not describe data collection, processing
        methodology or technical validation.
  \item \textbf{EC6:} Studies addressing exclusively non-clinical interactions
        or simulated contexts without demonstrated practical transferability.
\end{itemize}

\subsection{Information sources and search strategy}

The literature search was conducted across five electronic databases:
PubMed/MEDLINE (primary database), Google Scholar, EBSCO/CINAHL, Web of
Science, and Cochrane Library. Searches were performed independently by
three researchers, each applying multiple strategies covering the five
thematic axes of the review. General filters applied across all databases
were: publication period from January 2014 to July 2026, and articles
published in English, Spanish, or Portuguese.

Six search strategies were formulated. PubMed/MEDLINE was searched using
all six strategies: five thematic strings addressing specific conceptual axes
(clinical NLP and symptom extraction; named entity recognition and clinical
information extraction; clinician--patient communication and outcomes; vocal
biomarkers and speech analytics; NLP/speech-based systematic reviews), plus
one comprehensive unified string combining all major constructs of the
review. The remaining four databases (Google Scholar, EBSCO/CINAHL, Web of
Science, and Cochrane Library) were searched using the five thematic strings
only, with query syntax adapted to each platform's field codes and operators
while preserving the conceptual scope of each string. In an earlier search
iteration, a free full-text availability filter had been applied to
PubMed/MEDLINE owing to institutional access constraints; this restriction
was removed in the updated search to ensure comprehensive coverage.

The specific strings used in PubMed/MEDLINE syntax were as follows:

\begin{enumerate}
  \item \textbf{String 1 (NLP, Symptoms and EHR):}
        \begin{verbatim}
("natural language processing"[Title/Abstract])
AND
(symptoms[Title/Abstract] OR "patient-reported symptoms"[Title/Abstract])
AND
("electronic health records"[Title/Abstract] OR "patient-centered"[Title/Abstract])
        \end{verbatim}

  \item \textbf{String 2 (Clinical Information Extraction and NER):}
        \begin{verbatim}
("clinical information extraction"[Title/Abstract]
OR "medical information extraction"[Title/Abstract]
OR "named entity recognition"[Title/Abstract])
AND
("medical free text"[Title/Abstract]
OR "review"[Title/Abstract]
OR "patient-generated health data"[Title/Abstract])
        \end{verbatim}

  \item \textbf{String 3 (Clinician-Patient Communication and Outcomes):}
        \begin{verbatim}
("patient-physician communication"[Title/Abstract]
OR "doctor-patient communication"[Title/Abstract]
OR "patient-clinician relationship"[Title/Abstract])
AND
("intercultural settings"[Title/Abstract]
OR miscommunication[Title/Abstract]
OR "healthcare outcomes"[Title/Abstract]
OR "shared decision-making"[Title/Abstract]
OR "communication quality scale"[Title/Abstract])
        \end{verbatim}

  \item \textbf{String 4 (Vocal Biomarkers and Speech Analytics):}
        \begin{verbatim}
("vocal biomarkers"[Title/Abstract] OR "speech analytics"[Title/Abstract])
AND
("clinical practice"[Title/Abstract]
OR "voice for health"[Title/Abstract]
OR "SUS"[Title/Abstract])
        \end{verbatim}

  \item \textbf{String 5 (NLP, Speech and Reviews):}
        \begin{verbatim}
(("Applying natural language processing"[Title/Abstract])
  AND ("Systematic review"[Title/Abstract]))
OR
(("Identifying Topics"[Title/Abstract])
  AND ("Natural Language Processing Methods"[Title/Abstract]))
OR
(("Automatic speech recognition"[Title/Abstract])
  AND ("patient-clinician conversations"[Title/Abstract]))
        \end{verbatim}
\end{enumerate}

Additionally, a comprehensive search was conducted using the following
unified general string:

\begin{verbatim}
("speech analytics" OR "speech analysis" OR "speech recognition"
OR "automatic speech recognition" OR ASR OR "digital scribe"
OR "clinical transcription" OR "clinical documentation" OR "vocal biomarkers"
OR "healthcare conversations" OR "doctor-patient conversations"
OR "patient-doctor conversations")
AND
("natural language processing" OR NLP OR "large language model" OR LLM
OR "information extraction" OR "clinical information extraction"
OR "named entity recognition" OR NER OR "relation extraction" OR "text mining"
OR "conversation analysis" OR "topic modeling" OR "SOAP note" OR "SOAP notes")
AND
(healthcare OR clinical OR medicine OR medical OR physician OR clinician
OR doctor OR patient OR "electronic health record" OR EHR OR hospital
OR outpatient)
\end{verbatim}

To mitigate the loss of relevant articles and qualitatively enrich the corpus,
the electronic search was supplemented by the \textit{snowballing} method
(manual reference tracking or \textit{backward snowballing}). From the
systematic reviews and foundational articles located in the primary search,
their reference lists were manually analysed to identify and include other
studies that strictly meet the defined eligibility criteria, ensuring coverage
of key works of methodological or conceptual relevance to this review.

\subsection{Study selection}

Records retrieved from all five databases were exported and deduplicated
using Zotero reference management software, yielding a consolidated corpus
of 3,189 unique references. These records were imported into Rayyan
\citep{Ouzzani2016} for a second-pass deduplication and screening.

Titles and abstracts were screened independently by two reviewers, with
agreement assessed using Cohen's kappa coefficient ($\kappa = 0.78$);
disagreements were resolved by a third reviewer or the corresponding author.

To strengthen methodological rigour and assess the reliability of
LLM-assisted pre-screening as a decision-support tool, a locally-deployed
open-source large language model (Qwen2.5-Coder\,14B; Alibaba Cloud, 2024;
served via Ollama~v0.4) independently screened the same consolidated corpus
in parallel, configured with explicit inclusion and exclusion criteria
aligned with the review protocol \citep{Scherbakov2025}. The LLM produced a
structured decision (include/exclude/uncertain) with a confidence rating
and a one-sentence rationale for each record. Because per-record human
decisions were not centrally logged during the Rayyan screening, LLM
performance was instead validated by sensitivity analysis against the final
set of studies included in the synthesis: the LLM's recall was measured as
the proportion of the 24~included studies it correctly flagged as
``include'' or ``uncertain'' (records rated ``uncertain'' are, by protocol,
forwarded to human review rather than discarded). The LLM served exclusively
as a decision-support and validation tool: it did not participate in the
inclusion/exclusion decisions reported above, and all eligibility
determinations rest solely on human judgment.

Full texts of potentially eligible records were then retrieved and
assessed against the eligibility criteria. The selection process is
documented in the PRISMA-ScR flowchart (Fig.~\ref{fig:prisma}).

\subsection{Data extraction}

Data were extracted using a standardised form containing: author(s), year,
country, study design, clinical context, population characteristics,
technology used (ASR model, NLP framework, LLM), extracted clinical variables,
reported outcomes, ethical framework and limitations.

\subsection{Synthesis}

A narrative synthesis was performed. Findings are organised around the four
thematic axes defined in the eligibility criteria:
(A)~Clinical NLP and biomedical entity extraction;
(B)~automated clinical documentation and ASR;
(C)~physician-patient communication; and
(D)~vocal biomarkers and multimodal analysis.

%% ================================================================
\section{Results}
\label{sec:resultados}

% ---------------------------------------------------------------
\subsection{Study selection}
% ---------------------------------------------------------------

Searches across the five electronic databases (PubMed/MEDLINE, Scopus, Web
of Science, EBSCO/CINAHL, and Cochrane Library) retrieved 5,621~records
(Web of Science: 1,719; Scopus: 1,580; PubMed/MEDLINE: 1,059; EBSCO/CINAHL:
586; Cochrane Library: 399; Google Scholar: 278). After removing
2,432~duplicates, 3,189~unique records remained. These were screened by
title by two independent reviewers, resulting in the exclusion of
2,680~records (EC1: 1,120; EC2: 584; EC3: 134; EC4: 281; EC5: 358; EC6: 203)
and 509~records advancing to abstract screening. Abstract screening excluded
a further 474~records (EC1: 183; EC2: 116; EC3: 29; EC4: 43; EC5: 61;
EC6: 42), resulting in 35~articles selected for full-text assessment.
The 7~articles identified by \textit{snowballing} (backward reference
tracking of included systematic reviews) were added to the full-text
assessment stage independently of the electronic records. Of the 42~articles
assessed in full text, 18~were excluded for not meeting the eligibility
criteria (all from the electronic search; none from \textit{snowballing}).
In total, 24~articles were included in the synthesis
(17~from the electronic search + 7~from \textit{snowballing}). The
PRISMA-ScR flowchart is presented in Fig.~\ref{fig:prisma}.

The LLM-assisted pre-screening exercise correctly flagged 20 of the 24
included studies as ``include'' (83\% direct sensitivity); a further
3~studies (12.5\%) were rated ``uncertain'' and would, by protocol, have
been forwarded to human review rather than discarded, yielding an effective
recall of 95.8\% (23/24). The single false negative -- a study on
adherence metrics in orthopaedic care with no explicit NLP/ASR component in
its title -- was eligible under the broader clause of IC1 (studies that
measure the foundations of professional-patient communication) but fell
outside the narrower technology-focused exclusion criteria used to
configure the LLM. This recall is comparable to sensitivity figures
reported for LLM-assisted screening in the methodological literature
\citep{Issaiy2024}, and supports the use of LLM-assisted pre-screening as a
complementary, time-efficient decision-support tool rather than a
replacement for human screening.

\begin{figure}[h]
  \centering
  \resizebox{\linewidth}{!}{%% prisma_flow.tex — PRISMA-ScR Flowchart
%% Included via %% prisma_flow.tex — PRISMA-ScR Flowchart
%% Included via %% prisma_flow.tex — PRISMA-ScR Flowchart
%% Included via \input{prisma_flow} in the figure environment of main.tex
%% Exact counts: 5,621 identified across 5 databases, 2,432 duplicates removed,
%% 3,189 screened by title (2,680 excluded), 509 screened by abstract
%% (474 excluded), 35 full text (databases) + 7 snowballing, 18 excluded,
%% 24 included (17 databases + 7 snowballing).

\usetikzlibrary{positioning, calc, arrows.meta}

\begin{tikzpicture}[
  %% main flow boxes
  mbox/.style={
    rectangle, draw=black, thick, rounded corners=2pt,
    minimum width=6.8cm, text width=6.6cm,
    minimum height=1.3cm, align=center,
    fill=white, font=\small
  },
  %% exclusion boxes (right margin)
  ebox/.style={
    rectangle, draw=black, thick, rounded corners=2pt,
    minimum width=5.0cm, text width=4.8cm,
    minimum height=1.3cm, align=flush left,
    fill=gray!8, font=\small
  },
  %% arrows
  arr/.style={-{Stealth[length=5pt,width=4pt]}, thick},
  %% phase labels
  phase/.style={font=\bfseries\scriptsize, rotate=90, anchor=center}
]

%% ================================================================
%% PHASE 1 — IDENTIFICATION
%% ================================================================
\node[mbox, minimum width=6.0cm, text width=5.8cm] (dbbox) at (-1.8, 0) {
  \textbf{Electronic databases}\\[2pt]
  $n = 5{,}621$\\[2pt]
  {\scriptsize Web of Science: 1,719\\
    Scopus: 1,580\\
    PubMed/MEDLINE: 1,059\\
    EBSCO/CINAHL: 586\\
    Cochrane Library: 399\\
    Google Scholar: 278}
};

\node[mbox, minimum width=5.0cm, text width=4.8cm,
      right=1.5cm of dbbox] (otherbox) {
  \textbf{Other sources (Snowballing)}\\[2pt]
  $n = 7$\\[2pt]
  {\scriptsize (Manual backward reference\\
    tracking)}
};

%% ================================================================
%% PHASE 2 — SCREENING
%% ================================================================
\coordinate (midid) at ($(dbbox.south)!0.5!(otherbox.south)$);

\node[mbox, below=2.6cm of midid, anchor=north] (dedup) {
  \textbf{Database records after deduplication}\\[2pt]
  $n = 3{,}189$\quad{\scriptsize (5,621 $-$ 2,432 duplicates)}
};

%% arrow from databases to deduplication
\draw[arr] (dbbox.south) -- ++(0,-0.8cm) -| (dedup.north);

\node[mbox, below=1.2cm of dedup] (screen1) {
  \textbf{Title screening}\\[2pt]
  $n = 3{,}189$
};

\draw[arr] (dedup) -- (screen1);

\node[ebox, right=2.2cm of screen1] (excl1) {
  \textbf{Excluded by title:} $n = 2{,}680$\\[3pt]
  {\scriptsize
    EC1 (No NLP/ASR of clinical interactions): 1,120\\
    EC2 (No interaction-quality analysis): 584\\
    EC3 (Voice-based diagnosis only): 134\\
    EC4 (Publication type): 281\\
    EC5 (No methodological description): 358\\
    EC6 (Simulated/non-clinical context): 203
  }
};

\draw[arr] (screen1.east) -- (excl1.west);

\node[mbox, below=1.2cm of screen1] (screen2) {
  \textbf{Abstract screening}\\[2pt]
  $n = 509$
};

\draw[arr] (screen1) -- (screen2);

\node[ebox, right=2.2cm of screen2] (excl2) {
  \textbf{Excluded by abstract:} $n = 474$\\[3pt]
  {\scriptsize
    EC1 (No NLP/ASR of clinical interactions): 183\\
    EC2 (No interaction-quality analysis): 116\\
    EC3 (Voice-based diagnosis only): 29\\
    EC4 (Publication type): 43\\
    EC5 (No methodological description): 61\\
    EC6 (Simulated/non-clinical context): 42
  }
};

\draw[arr] (screen2.east) -- (excl2.west);

%% ================================================================
%% PHASE 3 — ELIGIBILITY
%% ================================================================
\node[mbox, below=1.2cm of screen2] (eligib) {
  \textbf{Full text assessed}\\[2pt]
  $n = 35$ (databases) $+$ $7$ (snowballing)
};

\draw[arr] (screen2) -- (eligib);

%% snowballing enters directly at eligibility (bypasses screening)
\draw[arr] (otherbox.south) -- ++(0,-0.8cm) -| ([xshift=1.1cm]eligib.east) -- (eligib.east);

\node[ebox, right=2.2cm of eligib] (excl3) {
  \textbf{Excluded in full text:} $n = 18$ (databases); $0$ (snowballing)\\[3pt]
  {\scriptsize
    Assessed against the full eligibility criteria (EC1--EC6)
  }
};

\draw[arr] (eligib.east) -- (excl3.west);

%% ================================================================
%% PHASE 4 — INCLUDED
%% ================================================================
\node[mbox, below=1.2cm of eligib] (included) {
  \textbf{Studies included in synthesis}\\[2pt]
  $n = 24$\\[2pt]
  {\scriptsize (17 databases + 7 snowballing)}
};

\draw[arr] (eligib) -- (included);

%% ================================================================
%% PHASE LABELS (left margin)
%% ================================================================
\node[phase] at ($(dbbox.west) + (-1.1cm, 0)$)      {IDENTIFICATION};
\node[phase] at ($(dedup.west)  + (-1.1cm, 0)$)      {SCREENING};
\node[phase] at ($(screen1.west) + (-1.1cm, 0)$)      {};   %% screening continues
\node[phase] at ($(screen2.west) + (-1.1cm, 0)$)      {};   %% screening continues
\node[phase] at ($(eligib.west) + (-1.1cm, 0)$)      {ELIGIBILITY};
\node[phase] at ($(included.west) + (-1.1cm, 0)$)    {INCLUDED};

%% ================================================================
%% FOOTNOTE
%% ================================================================
\node[font=\scriptsize, anchor=north west, text width=12cm]
  at ($(included.south west) + (-1.5cm, -0.5cm)$) {%
  Exact counts from the expanded multi-database search (2026).
  Snowballing assessed with the same eligibility criteria.
  $\kappa = 0.78$ (inter-reviewer agreement; Cohen's kappa).
};

\end{tikzpicture}
 in the figure environment of main.tex
%% Exact counts: 5,621 identified across 5 databases, 2,432 duplicates removed,
%% 3,189 screened by title (2,680 excluded), 509 screened by abstract
%% (474 excluded), 35 full text (databases) + 7 snowballing, 18 excluded,
%% 24 included (17 databases + 7 snowballing).

\usetikzlibrary{positioning, calc, arrows.meta}

\begin{tikzpicture}[
  %% main flow boxes
  mbox/.style={
    rectangle, draw=black, thick, rounded corners=2pt,
    minimum width=6.8cm, text width=6.6cm,
    minimum height=1.3cm, align=center,
    fill=white, font=\small
  },
  %% exclusion boxes (right margin)
  ebox/.style={
    rectangle, draw=black, thick, rounded corners=2pt,
    minimum width=5.0cm, text width=4.8cm,
    minimum height=1.3cm, align=flush left,
    fill=gray!8, font=\small
  },
  %% arrows
  arr/.style={-{Stealth[length=5pt,width=4pt]}, thick},
  %% phase labels
  phase/.style={font=\bfseries\scriptsize, rotate=90, anchor=center}
]

%% ================================================================
%% PHASE 1 — IDENTIFICATION
%% ================================================================
\node[mbox, minimum width=6.0cm, text width=5.8cm] (dbbox) at (-1.8, 0) {
  \textbf{Electronic databases}\\[2pt]
  $n = 5{,}621$\\[2pt]
  {\scriptsize Web of Science: 1,719\\
    Scopus: 1,580\\
    PubMed/MEDLINE: 1,059\\
    EBSCO/CINAHL: 586\\
    Cochrane Library: 399\\
    Google Scholar: 278}
};

\node[mbox, minimum width=5.0cm, text width=4.8cm,
      right=1.5cm of dbbox] (otherbox) {
  \textbf{Other sources (Snowballing)}\\[2pt]
  $n = 7$\\[2pt]
  {\scriptsize (Manual backward reference\\
    tracking)}
};

%% ================================================================
%% PHASE 2 — SCREENING
%% ================================================================
\coordinate (midid) at ($(dbbox.south)!0.5!(otherbox.south)$);

\node[mbox, below=2.6cm of midid, anchor=north] (dedup) {
  \textbf{Database records after deduplication}\\[2pt]
  $n = 3{,}189$\quad{\scriptsize (5,621 $-$ 2,432 duplicates)}
};

%% arrow from databases to deduplication
\draw[arr] (dbbox.south) -- ++(0,-0.8cm) -| (dedup.north);

\node[mbox, below=1.2cm of dedup] (screen1) {
  \textbf{Title screening}\\[2pt]
  $n = 3{,}189$
};

\draw[arr] (dedup) -- (screen1);

\node[ebox, right=2.2cm of screen1] (excl1) {
  \textbf{Excluded by title:} $n = 2{,}680$\\[3pt]
  {\scriptsize
    EC1 (No NLP/ASR of clinical interactions): 1,120\\
    EC2 (No interaction-quality analysis): 584\\
    EC3 (Voice-based diagnosis only): 134\\
    EC4 (Publication type): 281\\
    EC5 (No methodological description): 358\\
    EC6 (Simulated/non-clinical context): 203
  }
};

\draw[arr] (screen1.east) -- (excl1.west);

\node[mbox, below=1.2cm of screen1] (screen2) {
  \textbf{Abstract screening}\\[2pt]
  $n = 509$
};

\draw[arr] (screen1) -- (screen2);

\node[ebox, right=2.2cm of screen2] (excl2) {
  \textbf{Excluded by abstract:} $n = 474$\\[3pt]
  {\scriptsize
    EC1 (No NLP/ASR of clinical interactions): 183\\
    EC2 (No interaction-quality analysis): 116\\
    EC3 (Voice-based diagnosis only): 29\\
    EC4 (Publication type): 43\\
    EC5 (No methodological description): 61\\
    EC6 (Simulated/non-clinical context): 42
  }
};

\draw[arr] (screen2.east) -- (excl2.west);

%% ================================================================
%% PHASE 3 — ELIGIBILITY
%% ================================================================
\node[mbox, below=1.2cm of screen2] (eligib) {
  \textbf{Full text assessed}\\[2pt]
  $n = 35$ (databases) $+$ $7$ (snowballing)
};

\draw[arr] (screen2) -- (eligib);

%% snowballing enters directly at eligibility (bypasses screening)
\draw[arr] (otherbox.south) -- ++(0,-0.8cm) -| ([xshift=1.1cm]eligib.east) -- (eligib.east);

\node[ebox, right=2.2cm of eligib] (excl3) {
  \textbf{Excluded in full text:} $n = 18$ (databases); $0$ (snowballing)\\[3pt]
  {\scriptsize
    Assessed against the full eligibility criteria (EC1--EC6)
  }
};

\draw[arr] (eligib.east) -- (excl3.west);

%% ================================================================
%% PHASE 4 — INCLUDED
%% ================================================================
\node[mbox, below=1.2cm of eligib] (included) {
  \textbf{Studies included in synthesis}\\[2pt]
  $n = 24$\\[2pt]
  {\scriptsize (17 databases + 7 snowballing)}
};

\draw[arr] (eligib) -- (included);

%% ================================================================
%% PHASE LABELS (left margin)
%% ================================================================
\node[phase] at ($(dbbox.west) + (-1.1cm, 0)$)      {IDENTIFICATION};
\node[phase] at ($(dedup.west)  + (-1.1cm, 0)$)      {SCREENING};
\node[phase] at ($(screen1.west) + (-1.1cm, 0)$)      {};   %% screening continues
\node[phase] at ($(screen2.west) + (-1.1cm, 0)$)      {};   %% screening continues
\node[phase] at ($(eligib.west) + (-1.1cm, 0)$)      {ELIGIBILITY};
\node[phase] at ($(included.west) + (-1.1cm, 0)$)    {INCLUDED};

%% ================================================================
%% FOOTNOTE
%% ================================================================
\node[font=\scriptsize, anchor=north west, text width=12cm]
  at ($(included.south west) + (-1.5cm, -0.5cm)$) {%
  Exact counts from the expanded multi-database search (2026).
  Snowballing assessed with the same eligibility criteria.
  $\kappa = 0.78$ (inter-reviewer agreement; Cohen's kappa).
};

\end{tikzpicture}
 in the figure environment of main.tex
%% Exact counts: 5,621 identified across 5 databases, 2,432 duplicates removed,
%% 3,189 screened by title (2,680 excluded), 509 screened by abstract
%% (474 excluded), 35 full text (databases) + 7 snowballing, 18 excluded,
%% 24 included (17 databases + 7 snowballing).

\usetikzlibrary{positioning, calc, arrows.meta}

\begin{tikzpicture}[
  %% main flow boxes
  mbox/.style={
    rectangle, draw=black, thick, rounded corners=2pt,
    minimum width=6.8cm, text width=6.6cm,
    minimum height=1.3cm, align=center,
    fill=white, font=\small
  },
  %% exclusion boxes (right margin)
  ebox/.style={
    rectangle, draw=black, thick, rounded corners=2pt,
    minimum width=5.0cm, text width=4.8cm,
    minimum height=1.3cm, align=flush left,
    fill=gray!8, font=\small
  },
  %% arrows
  arr/.style={-{Stealth[length=5pt,width=4pt]}, thick},
  %% phase labels
  phase/.style={font=\bfseries\scriptsize, rotate=90, anchor=center}
]

%% ================================================================
%% PHASE 1 — IDENTIFICATION
%% ================================================================
\node[mbox, minimum width=6.0cm, text width=5.8cm] (dbbox) at (-1.8, 0) {
  \textbf{Electronic databases}\\[2pt]
  $n = 5{,}621$\\[2pt]
  {\scriptsize Web of Science: 1,719\\
    Scopus: 1,580\\
    PubMed/MEDLINE: 1,059\\
    EBSCO/CINAHL: 586\\
    Cochrane Library: 399\\
    Google Scholar: 278}
};

\node[mbox, minimum width=5.0cm, text width=4.8cm,
      right=1.5cm of dbbox] (otherbox) {
  \textbf{Other sources (Snowballing)}\\[2pt]
  $n = 7$\\[2pt]
  {\scriptsize (Manual backward reference\\
    tracking)}
};

%% ================================================================
%% PHASE 2 — SCREENING
%% ================================================================
\coordinate (midid) at ($(dbbox.south)!0.5!(otherbox.south)$);

\node[mbox, below=2.6cm of midid, anchor=north] (dedup) {
  \textbf{Database records after deduplication}\\[2pt]
  $n = 3{,}189$\quad{\scriptsize (5,621 $-$ 2,432 duplicates)}
};

%% arrow from databases to deduplication
\draw[arr] (dbbox.south) -- ++(0,-0.8cm) -| (dedup.north);

\node[mbox, below=1.2cm of dedup] (screen1) {
  \textbf{Title screening}\\[2pt]
  $n = 3{,}189$
};

\draw[arr] (dedup) -- (screen1);

\node[ebox, right=2.2cm of screen1] (excl1) {
  \textbf{Excluded by title:} $n = 2{,}680$\\[3pt]
  {\scriptsize
    EC1 (No NLP/ASR of clinical interactions): 1,120\\
    EC2 (No interaction-quality analysis): 584\\
    EC3 (Voice-based diagnosis only): 134\\
    EC4 (Publication type): 281\\
    EC5 (No methodological description): 358\\
    EC6 (Simulated/non-clinical context): 203
  }
};

\draw[arr] (screen1.east) -- (excl1.west);

\node[mbox, below=1.2cm of screen1] (screen2) {
  \textbf{Abstract screening}\\[2pt]
  $n = 509$
};

\draw[arr] (screen1) -- (screen2);

\node[ebox, right=2.2cm of screen2] (excl2) {
  \textbf{Excluded by abstract:} $n = 474$\\[3pt]
  {\scriptsize
    EC1 (No NLP/ASR of clinical interactions): 183\\
    EC2 (No interaction-quality analysis): 116\\
    EC3 (Voice-based diagnosis only): 29\\
    EC4 (Publication type): 43\\
    EC5 (No methodological description): 61\\
    EC6 (Simulated/non-clinical context): 42
  }
};

\draw[arr] (screen2.east) -- (excl2.west);

%% ================================================================
%% PHASE 3 — ELIGIBILITY
%% ================================================================
\node[mbox, below=1.2cm of screen2] (eligib) {
  \textbf{Full text assessed}\\[2pt]
  $n = 35$ (databases) $+$ $7$ (snowballing)
};

\draw[arr] (screen2) -- (eligib);

%% snowballing enters directly at eligibility (bypasses screening)
\draw[arr] (otherbox.south) -- ++(0,-0.8cm) -| ([xshift=1.1cm]eligib.east) -- (eligib.east);

\node[ebox, right=2.2cm of eligib] (excl3) {
  \textbf{Excluded in full text:} $n = 18$ (databases); $0$ (snowballing)\\[3pt]
  {\scriptsize
    Assessed against the full eligibility criteria (EC1--EC6)
  }
};

\draw[arr] (eligib.east) -- (excl3.west);

%% ================================================================
%% PHASE 4 — INCLUDED
%% ================================================================
\node[mbox, below=1.2cm of eligib] (included) {
  \textbf{Studies included in synthesis}\\[2pt]
  $n = 24$\\[2pt]
  {\scriptsize (17 databases + 7 snowballing)}
};

\draw[arr] (eligib) -- (included);

%% ================================================================
%% PHASE LABELS (left margin)
%% ================================================================
\node[phase] at ($(dbbox.west) + (-1.1cm, 0)$)      {IDENTIFICATION};
\node[phase] at ($(dedup.west)  + (-1.1cm, 0)$)      {SCREENING};
\node[phase] at ($(screen1.west) + (-1.1cm, 0)$)      {};   %% screening continues
\node[phase] at ($(screen2.west) + (-1.1cm, 0)$)      {};   %% screening continues
\node[phase] at ($(eligib.west) + (-1.1cm, 0)$)      {ELIGIBILITY};
\node[phase] at ($(included.west) + (-1.1cm, 0)$)    {INCLUDED};

%% ================================================================
%% FOOTNOTE
%% ================================================================
\node[font=\scriptsize, anchor=north west, text width=12cm]
  at ($(included.south west) + (-1.5cm, -0.5cm)$) {%
  Exact counts from the expanded multi-database search (2026).
  Snowballing assessed with the same eligibility criteria.
  $\kappa = 0.78$ (inter-reviewer agreement; Cohen's kappa).
};

\end{tikzpicture}
}
  \caption{PRISMA-ScR flowchart of the study selection process.}
  \label{fig:prisma}
\end{figure}

% ---------------------------------------------------------------
\subsection{Characteristics of included studies}
% ---------------------------------------------------------------

The 24~studies in the corpus span the period 2014--2025
(Table~\ref{tab:estudos}).

\begin{longtable}{p{0.65cm} p{2.4cm} p{1.55cm} p{0.55cm} p{2.2cm} p{2.35cm} p{1.4cm}}
\caption{Characteristics of studies included in the synthesis.}
\label{tab:estudos}\\
\toprule
\textbf{Code} & \textbf{Author, year} & \textbf{Type} &
\textbf{Axis} & \textbf{Technology} & \textbf{Main outcome} & \textbf{Country} \\
\midrule
\endfirsthead
\multicolumn{7}{l}{\small\textit{(continued)}}\\
\toprule
\textbf{Code} & \textbf{Author, year} & \textbf{Type} &
\textbf{Axis} & \textbf{Technology} & \textbf{Main outcome} & \textbf{Country} \\
\midrule
\endhead
\bottomrule
\endfoot
%% ---- Axis A: Clinical NLP and biomedical entity extraction ----
S01 & Bazoge et al., 2023    & Syst.\ rev.  & A & Transformers, BERT         & NER in CDWs; phenotyping               & France        \\
S02 & Navarro et al., 2023   & Syst.\ rev.  & A & BERT, BioBERT              & Clinical NER (problems, treatments)    & Australia     \\
S03 & Watabe et al., 2025    & Empirical    & A & BERT-CRF                   & F1\,>\,0.8; reported symptoms          & Japan         \\
S04 & Sezgin et al., 2023    & Empirical    & A & NER + SNOMED CT (zero-shot) & Symptom identification in PGHD        & USA           \\
S05 & Koleck et al., 2019    & Syst.\ rev.  & A & NLP, classifiers           & Symptom extraction from EHR            & USA           \\
S06 & Wang et al., 2018      & Lit.\ rev.   & A & cTAKES, MedLEE, ML         & 263 clinical extraction systems        & USA           \\
%% ---- Axis B: Automated clinical documentation and ASR ----
S07 & Krishna et al., 2021   & Preprint\textsuperscript{‡} & B & CLUSTER2SENT    & SOAP notes generated from consultations & USA          \\
S08 & Klusty et al., 2025    & Empirical    & B & Whisper + PyAnnote         & On-premise transcription; privacy      & USA           \\
S09 & Abbasian et al., 2024  & Framework    & B & Generative AI, RAG         & Groundedness, safety, empathy          & USA           \\
S10 & Aleksić et al., 2017  & Empirical & B & EHR summarisation & Dashboard for chronic disease & Serbia  \\
S11 & Tran et al., 2023      & Empirical    & B & General vs.\ specialist ASR & WER; errors in short pronouns         & USA           \\
S12 & Schloss \& Konam, 2020 & Preprint\textsuperscript{‡} & B & Hierarchical Bi-LSTM & Utterance classification → SOAP  & USA          \\
S13 & van Buchem et al., 2021 & Scoping rev. & B & Digital scribes (review) & Usability gaps in real clinical use   & Netherlands   \\
S14 & Falcetta et al., 2023  & Syst.\ rev.  & B & Scribes (review)           & Mapping of 22 scribes; no prospective validation & Brazil \\
S15 & Ng et al., 2025        & Syst.\ rev.  & B & Clinical ASR (review)      & WER ranged from 8.7\% to $>$50\% by context & Singapore \\
%% ---- Axis C: Physician-patient communication ----
S16 & Shao et al., 2025      & Empirical    & C & DPCQ scale (14 items)      & Physician-patient communication quality & China       \\
S17 & Muscat et al., 2020    & Framework    & C & Shared decision-making     & Health literacy; engagement            & Australia   \\
S18 & Kelley et al., 2014    & Meta-analysis & C & --                        & Clinical relationship → objective outcomes & USA     \\
S19 & Rucinski et al., 2022  & Syst.\ rev.  & C & Adherence metrics          & Bilateral adherence in orthopedics     & USA           \\
S20 & McCabe \& Healey, 2018 & Empirical    & C & Conversation analysis      & Self-repair as adherence predictor     & UK            \\
S21 & Alkhamees \& Alasqah, 2023 & Integr.\ rev. & C & --               & Intercultural barriers in communication & Saudi A.    \\
%% ---- Axis D: Vocal biomarkers and multimodal analysis ----
S22 & Fagherazzi et al., 2021 & Framework   & D & Vocal pipeline (acoustics+NLP) & Vocal biomarkers of health and emotion & Luxembourg \\
S23 & Albert et al., 2024    & Scoping rev.\textsuperscript{‡} & D & AI, Speech, NLP & Clinician-patient communication via AI & France  \\
S24 & Imel et al., 2022      & Report       & D & Sequential neural networks & Emotional valence; human rater agreement & USA         \\
\midrule
\multicolumn{7}{l}{\footnotesize\parbox{13.5cm}{%
  \textsuperscript{‡}Preprint/not peer-reviewed (arXiv/medRxiv) --- included as they meet IC1--IC4; preprint acceptance per Elsevier/AIM policy.}} \\
\end{longtable}

% ---------------------------------------------------------------
\subsection{Axis A —Clinical NLP and biomedical entity extraction}
\label{subsec:eixoA}
% ---------------------------------------------------------------

The studies in Axis~A map the evolution of NLP approaches applied to clinical
information extraction, documenting the transition from rule-based systems to
deep learning architectures and \textit{transformer} models.

The historical landscape of the field is described by \citet{Wang2018}, whose
review of 263~studies evidences the predominance of symbolic systems such as
cTAKES and MedLEE in the processing of electronic health records. The authors
demonstrate how adherence to linguistically rule-based approaches was the norm
until the arrival of machine learning techniques. The consolidation of the
transition to deep learning models is documented by \citet{Bazoge2023}, in a
systematic review on the use of NLP in clinical data warehouses (CDWs): the
authors evidence the migration from symbolic methods to \textit{transformers},
identifying their applicability in phenotyping and symptom management tasks
from secondary data.

In the specific domain of biomedical Named Entity Recognition (NER),
\citet{Navarro2023} systematically reviewed the literature on NER and relation
extraction in free-text medical records and concluded that models such as BERT
lead in extracting the categories of ``problems'' and ``treatments''. The
authors point out, however, the scarcity of implementations in real clinical
settings as a persistent limitation of the field.

The ability of \textit{transformer} models to capture patient colloquial
language was demonstrated by \citet{Watabe2025}, who developed and validated a
BERT-CRF model for extracting symptoms reported in everyday language, achieving
an F1-score above~0.8. The study highlights the need for \textit{fine-tuning}
with authentic narrative data to capture slang and linguistic nuances present
in patient speech. Complementarily, \citet{Sezgin2023} demonstrated the
feasibility of a \textit{zero-shot} approach combining NER and the SNOMED~CT
ontology to identify symptoms in patient-generated health data (PGHD), without
specific retraining, validating the possibility of pre-trained models
translating lay language into structured clinical metrics.

At the intersection of NLP and electronic health records, \citet{Koleck2019}
synthesised the efficacy of classification algorithms for extracting subjective
symptoms -- such as pain and mood changes -- from free-text EHR narratives,
evidencing the technical feasibility of automated extraction even for
subjective dimensions.

Taken together, the studies in Axis~A confirm the technical maturity of
clinical NLP for biomedical entity extraction, with BERT and its variants
established as the state of the art. The cross-cutting gap identified --
highlighted by \citet{Navarro2023} -- is the distance between performance in
laboratory settings and deployment in real clinical systems, where data
privacy, linguistic variability and integration with care workflows remain
unsolved challenges.

% ---------------------------------------------------------------
\subsection{Axis B —Automated clinical documentation and ASR}
\label{subsec:eixoB}
% ---------------------------------------------------------------

The studies in Axis~B cover the full cycle of automated clinical
documentation: from speech recognition to the generation of structured
records, including the technical challenges and ethical requirements for
hospital deployment.

The foundations of automated documentation are established by
\citet{Aleksic2017}, whose electronic health record summarisation model
demonstrates the feasibility of dashboards capable of identifying patients
requiring specialist attention, validating readable synthesis as a strategy
to reduce clinician cognitive load.

In the specific domain of SOAP note generation, two studies address
complementary strategies. \citet{Schloss2020} used hierarchical Bi-LSTM
networks to classify consultation utterances into SOAP sections with high
accuracy, proposing modular strategies to handle data ``noise'' in automatic
transcriptions. \citet{Krishna2021}, in turn, proposed the CLUSTER2SENT
model, based on modular summarisation, to generate structured notes from
dense physician-patient conversations, facilitating the retrieval of verbal
evidence by the clinician in the analytical output.

With regard to automatic transcription, \citet{Klusty2025} described a system
based on OpenAI Whisper and PyAnnote operating on closed (\textit{on-premise})
servers, simultaneously resolving the speaker diarisation challenge and the
privacy requirements for hospital deployment. Complementarily,
\citet{Tran2023} demonstrated that errors in short pronouns (such as
\enquote{I} and \enquote{you}) are more critical for the validity of clinical
indicators than errors in complex medical terms, highlighting the need for
accuracy in casual language to ensure the quality of extracted data.

Generative AI-based approaches bring a distinct set of metric requirements.
\citet{Abbasian2024} propose safety, factual validity
(\textit{groundedness}) and empathy metrics as mandatory indicators for
evaluating health conversations generated by language models, grounding the
introduction of Retrieval-Augmented Generation (RAG) as a mechanism to anchor
model responses in facts extracted from the consultation and mitigate
algorithmic hallucinations.

The state of the field is synthesised by \citet{vanBuchem2021}, in a scoping
review on the deployment of digital scribes in clinical practice: the authors
point to the absence of prospective usability and real clinical utility studies
as the central gap in the field, with most initiatives still limited to
proof-of-concept without prospective clinical validation.

Taken together, the studies in Axis~B reveal a field in rapid technical
evolution (with ASR, diarisation and SOAP note generation \textit{pipelines}
progressively more robust), but with prospective clinical validation still
scarce \citep{vanBuchem2021}. Most approaches direct processing towards
documentation automation, subordinating the analysis of clinical interaction
quality to the production of the structured record.

% ---------------------------------------------------------------
\subsection{Axis C —Physician-patient communication}
\label{subsec:eixoC}
% ---------------------------------------------------------------

The studies in Axis~C establish the empirical basis for professional-patient
communication quality as a determinant of clinical outcomes, identifying the
measurable dimensions of the therapeutic interaction and the barriers that
compromise its effectiveness.

The impact of the clinician-patient relationship on objective health outcomes
is demonstrated by the meta-analysis of \citet{Kelley2014}, which confirms
the significant effect of the emotional and cognitive dimensions of
communication on measurable clinical results, providing the theoretical
foundation for investing in automatic quality metrics of the therapeutic
relationship.

In the domain of structured measurement, \citet{Shao2025} developed and
validated the 14-item DPCQ (\textit{Doctor-Patient Communication Quality})
scale, focused on emotional support and health education. The semantic
categories proposed by the authors (emotional interaction versus behavioural
guidance) are directly applicable as targets for speech-derived automatic
metrics.

The dimension of therapeutic adherence in the orthopedic context is addressed
by \citet{Rucinski2022}, whose systematic review demonstrates that current
metrics fail by ignoring the barriers imposed by the healthcare team and the
system, not only by the patient. The authors argue for the creation of
auditable data capable of mapping bilateral failures in orthopedic
rehabilitation, an argument that reinforces the need for instruments to
analyse the interaction, not just the isolated behaviour of the patient.

At the microscopic level of speech, \citet{McCabe2018} demonstrated that the
use of \enquote{self-repair} by physicians (i.e., the spontaneous correction
of utterances during the consultation) predicts superior treatment adherence.
The finding validates the analysis of low-level linguistic dynamics as
predictors of clinical outcomes, demonstrating that automatically identifiable
communicational markers can carry prognostic value.

The role of health literacy in effective communication is analysed by
\citet{Muscat2020}, who demonstrate that shared decision-making requires the
active building of patient skills through effective verbal instruction, in
contrast with the passive informational model. Speech monitoring in the
clinical context must therefore capture indicators of communicational literacy
success, not merely unilateral information transfer.

Linguistic and cultural diversity contexts introduce an additional layer of
complexity. \citet{Alkhamees2023} mapped how linguistic disparities and
authoritarian consultation styles generate patient inhibition and compromise
the quality of clinical decisions. In the diverse context of SUS (Brazil's
National Health System), these barriers often remain invisible to managers,
but are potentially identifiable through automated interaction analysis.

The clinical propedeutics systematised by \citet{Porto2019} organises the
anamnesis as a structured sequence of inquiry elements, defining the
propedeutic characteristics of the complaint -- location, intensity, quality,
radiation, triggering and relieving factors, duration and evolution -- as
mandatory semantic targets of diagnostic reasoning. This structure constitutes
the clinical grammar of verbal interaction: the set of information that the
clinician must elicit and the patient must verbalise for the anamnesis to
fulfil its diagnostic function. None of the studies in this review's corpus
crosses the propedeutic model with automatic analysis of verbal interaction:
the studies in Axis~C measure communication quality in its relational and
cognitive dimensions, but do not verify whether the propedeutic characteristics
were effectively verbalised, understood and recorded during the consultation.

Taken together, the studies in Axis~C provide the map of what must be measured
in clinical interaction: emotional and cognitive dimensions
\citep{Kelley2014,Shao2025}, linguistic adherence markers
\citep{McCabe2018,Rucinski2022}, health literacy competencies
\citep{Muscat2020} and intercultural barriers \citep{Alkhamees2023}. The
emerging gap is twofold: no study operationalises these dimensions as
automatically speech-derived variables (technological gap), and no study
verifies whether the propedeutic characteristics of the anamnesis
\citep{Porto2019} were elicited with adequate communicational quality during
real interaction (clinical-methodological gap). Measurement remains dependent
on self-report scales or direct observation, without integration with
acoustic-semantic analysis \textit{pipelines} guided by the clinical structure
of the anamnesis.

% ---------------------------------------------------------------
\subsection{Axis D —Vocal biomarkers and multimodal analysis}
\label{subsec:eixoD}
% ---------------------------------------------------------------

The studies in Axis~D explore the acoustic and prosodic dimension of clinical
communication, extending the analysis beyond the semantic content of the
transcribed text to the vocal signals that accompany speech.

\citet{Fagherazzi2021} describe the \textit{pipeline} for capturing linguistic
and acoustic features -- prosody, vocal quality and rhythm -- that indicate
health states and emotional states of the speaker. The authors demonstrate
that signals derived directly from the voice can function as clinical
biomarkers, grounding the idea that the speech of both patient and clinician,
beyond its semantic content, carries usable analytical information
independently of textual transcription.

In the domain of real clinical consultation analysis, \citet{Imel2022} mapped
emotional valences from textual transcripts of physician-patient dialogues
using sequential neural networks, achieving agreement with trained human
raters. Although the study demonstrates that emotional valence markers of
verbal interaction can be automatically identified at utterance level with
reliability comparable to human assessment (correlation of 0.60), the authors
note that the model was trained and evaluated exclusively on the textual
content of the dialogues (transcription diaries), without directly processing
the audio signal or the acoustic layer.

The unresolved frontier of Axis~D is identified by \citet{Albert2024}, in a
scoping review on AI, speech and NLP methods for clinician-patient
communication assessment: the authors highlight the scarcity of studies on
non-verbal language, voice tones and silences. Most existing approaches focus
on transcription and semantic analysis, neglecting the acoustic-prosodic layer
as an independent object of analysis. Silences, pauses and voice modulations --
therapeutic and relational components of the anamnesis -- remain without a
consolidated method for automatic extraction and interpretation.

Taken together, the studies in Axis~D reveal a field in an early stage of
development: the technical foundations of voice processing exist
\citep{Fagherazzi2021} and there is feasibility in mapping emotional markers
based on textual transcripts of real interactions \citep{Imel2022}, but the
multimodal integration that truly unites the acoustic aspects of voice with
the semantic aspects of text in clinical communication remains incipient
\citep{Albert2024}. This gap acquires precise clinical contour when confronted
with the propedeutics of \citet{Porto2019}: if the propedeutic characteristics
of the complaint have acoustic-prosodic correlates in speech (vocal tension for
pain, hesitation and pauses for doubt, prosodic disengagement for barriers),
no study in the corpus guides the extraction of vocal biomarkers by the
clinical grammar of the anamnesis. The integration between the propedeutic
model and the acoustic-semantic \textit{pipeline} represents the most
unexplored methodological frontier of this review.

%% ================================================================
\section{Discussion}
\label{sec:discussao}


\subsection*{Two parallel axes without systematic convergence}

The mapped studies reveal two lines of development advancing in parallel
without systematic convergence in the literature.
The first line, formed by clinical NLP studies (Axis~A) and automated
documentation studies (Axis~B), treats speech as input for documentation:
the clinical encounter is transcribed to generate SOAP notes, identify
biomedical entities or populate electronic health record fields
\citep{Bazoge2023,vanBuchem2021,Falcetta2023,Ng2025}.
The second line, formed by physician-patient communication studies (Axis~C)
and vocal biomarker studies (Axis~D), establishes what communication means
clinically and demonstrates that acoustic signals carry diagnostic information
\citep{Kelley2014,Fagherazzi2021,Albert2024}.
In this tradition, communication quality is operationalised by self-report
instruments or in-person observation, and vocal biomarkers are applied
predominantly to the detection of health states (Parkinson's disease,
depression, COVID-19), not to the communicational dynamics during the
consultation.

The methodological heterogeneity reflects this structural divide.
NLP studies such as \citet{Bazoge2023} and
\citet{Wang2018} employ entity extraction pipelines on retrospective
EHR text; ASR studies such as \citet{vanBuchem2021} and
\citet{Klusty2025} measure transcription accuracy and documentation time
reduction; communication studies such as \citet{Shao2025} and
\citet{Kelley2014} derive relational indicators through validated
questionnaires. These traditions share the locus of clinical interaction
but do not share measurement objects, making quantitative synthesis
unfeasible with the studies in the current corpus.

\subsection*{Confirmed gap: anamnesis as the object of analysis}

No study in the corpus operationalises the anamnesis as a structured clinical
event for the automatic extraction of speech-derived communicational and
behavioural markers.
The distinction is critical: most studies treat the clinical interaction as a
\emph{means} -- to produce documentation or detect pathologies -- not as an
\emph{object} of communicational analysis.
\citet{Krishna2021} and
\citet{Schloss2020} classify utterances by SOAP structure, not by the
relational quality of the encounter.
\citet{Abbasian2024} evaluate the efficacy of conversations generated by
generative AI, not the communicative behaviour of the healthcare professional
in real context.
\citet{Shao2025} and \citet{McCabe2018} describe relevant dimensions of
communicative quality -- clarity, empathy, comprehension barriers -- but
operationalise them through self-report instruments, not through automatically
extracted speech markers.

The review by \citet{Albert2024} is the one that comes closest to this
intersection: it maps SA and NLP methods for clinical communication assessment,
but confirms the scarcity of studies linking acoustic-prosodic markers to
measurable communicational outcomes, particularly in rehabilitation contexts.
\citet{Porto2019} offers a precise semiological framework for orthopedic
anamnesis, including symptom hierarchy, chronology and functional exploration,
which could guide the operational definition of these markers; however, no
study retrieved in this review integrates that framework with automated
speech analysis.

\subsection*{Implications for prospective studies}

The identified gap points to a set of requirements for future studies:
(i)~recording of real clinical encounters under ecological conditions;
(ii)~ASR with models fine-tuned to Brazilian Portuguese;
(iii)~extraction of communicational markers (expression of doubt, pain,
engagement, linguistic barriers) anchored in the semiological framework of
the anamnesis \citep{Porto2019}; and
(iv)~correlation of these markers with adherence outcomes and quality of care
in longitudinal follow-up.
The convergence of requirements identified across the four axes --
ecological recording conditions, ASR fine-tuned to Brazilian Portuguese,
semiologically anchored marker extraction, and longitudinal outcome
correlation -- defines the methodological frame that Phase 2 prospective
studies must address. A protocol meeting this frame has been proposed and
published for orthopedic rehabilitation in SUS [Author et al., 2026]. The
present review provides the evidence synthesis that validates the
scientific rationale of that protocol and operationalises the specific
gaps it must address, namely: the integration of ASR-derived transcription
with communicational marker extraction anchored in the propedeutic model
of \citet{Porto2019}, tested against adherence and quality-of-care
outcomes in longitudinal follow-up.

\subsection*{Limitations of this review}

This review presents three central limitations.
First, although the literature search was substantially expanded to cover
multiple indexed electronic databases (PubMed/MEDLINE, Scopus, Web of
Science, EBSCO/CINAHL, Cochrane Library, and Google Scholar) without a
free full-text filter, the final inclusion count remained at 17
database-derived studies owing to the strict application of eligibility
criteria during the screening and full-text reading phases.
Second, although the search included English, Spanish, and Portuguese, the
mapped literature remains predominantly in English, with under-representation
of studies conducted in the Brazilian health context, which may underestimate
local Speech Analytics initiatives applied to SUS.
Third, the heterogeneity of study types --- including systematic reviews,
empirical studies, preprints, and conceptual models --- prevents quantitative
synthesis or structured meta-analysis, limiting direct comparisons of
methodological quality between approaches.

\subsection*{Ethical considerations}

The application of Speech Analytics in clinical consultations raises ethical
questions that go beyond standard human research considerations.
Voice constitutes sensitive biometric data under the terms of the General
Personal Data Protection Law (LGPD, Law No.~13.709/2018), imposing
obligations of data minimisation, pseudonymisation, separation between
identified and analytical datasets, and restricted access control.
Research must be conducted in compliance with CNS Resolution
No.~466/2012 \citep{CNS466_2012} and CNS Resolution No.~510/2016
\citep{CNS510_2016}, and with Law No.~14.874/2024 \citep{Lei14874_2024},
which updates the regulatory framework for clinical research in Brazil.
The Free and Informed Consent Form (TCLE -- \textit{Termo de Consentimento
Livre e Esclarecido}) is mandatory for voice recording and cannot be waived
when speech itself constitutes the object of analysis.

From a technical standpoint, \citet{Klusty2025}
demonstrate that on-premise processing is feasible in a hospital environment
and represents the most appropriate solution for privacy, eliminating the
transmission of sensitive data to external servers.
This approach aligns with the five principles for responsible AI proposed by
\citet{Floridi2018} (beneficence, non-maleficence, autonomy, justice and
explicability), and with the international recommendations of the OECD
\citep{OECD2024} and the World Health Organization \citep{WHO2024} for
ethical governance of large-scale AI models in healthcare.
The research protocol must also explicitly prohibit the use of extracted
markers for autonomous clinical decision-making, preserving the
decision-making role of the healthcare professional and the informed autonomy
of the patient.

%% ================================================================
\section{Conclusion}
\label{sec:conclusao}

This scoping review mapped the state of the art in Speech Analytics and
clinical communication with the aim of establishing the scientific foundation
for prospective studies that integrate automated speech analysis with the
assessment of communicational quality in orthopedic rehabilitation contexts.

The four thematic axes evidence consolidated advances in clinical NLP and ASR:
biomedical entity extraction \textit{pipelines} achieve competitive performance
in electronic health records \citep{Bazoge2023,Koleck2019}, speech recognition
models such as Whisper \citep{Radford2023} enable the transcription of clinical
encounters in hospital settings \citep{Klusty2025,Ng2025}, and vocal
biomarkers demonstrate the capacity to capture physiological and emotional
states relevant to healthcare \citep{Fagherazzi2021,Albert2024}.
The physician-patient communication literature, in turn, consistently
establishes that the relational quality of the clinical encounter influences
therapeutic adherence and health outcomes \citep{Kelley2014,Rucinski2022}.

However, the review confirms a consistent and specific gap: few studies
operationalise clinical interactions as a source of speech-derived measurable
variables, particularly communicational and behavioural markers (expression of
doubt, pain, engagement, comprehension barriers) automatically extracted
during the anamnesis as a structured event.
Most initiatives treat speech as input for documentation or as a signal for
pathology detection, without anchoring the analysis in the semiological
framework of the anamnesis \citep{Porto2019}, which defines what, in speech,
is clinically significant from a communicational standpoint.

This gap points to the need for prospective studies that integrate Speech
Analytics with structured analysis of clinical communication for real-time
adherence monitoring and quality of care.
Such studies require, in addition to adequate technical infrastructure, a
robust ethical protocol that treats voice as sensitive biometric data, ensures
informed consent and preserves patient privacy through local data processing.
This review constitutes the consolidated evidence base for the prospective
empirical phase of this research programme, delineating the methodological
requirements that future studies must address and establishing the
scientific legitimacy of integrating \textit{Speech Analytics} with
structured clinical communication analysis in the context of SUS
(Brazil's National Health System).

%% ================================================================
\section*{Declarations}

\subsection*{Funding}
This study received no funding from funding agencies or entities from the
public, commercial or non-profit sectors.
Article processing charges (APC) are waived for the corresponding author's
institution under the Open Access Agreement signed with Elsevier for the
period 2026--2028 (\textit{Read \& Publish Agreement}).

\subsection*{Declaration of competing interest}
None of the authors has employment, consultancy, equity, honoraria, patents
or funding ties to entities with commercial or institutional interest in the
topic of this review.
The prospective study protocol cited in the text [Author et al., 2026] is
co-authored by members of the authorship team; this relationship will be
fully disclosed in the author list upon acceptance and does not constitute
a conflict of interest under ICMJE guidelines.

\subsection*{Ethical approval}
Not applicable. This study consists of a scoping review based exclusively on
published and publicly accessible scientific literature. No primary data were
collected, no participants were recruited, and no intervention was performed
on human subjects.
Under the terms of CNS Resolution No.~510/2016 \citep{CNS510_2016}, studies
using public domain information are exempt from registration and assessment
by the CEP/CONEP system.

\subsection*{Data availability}
The data supporting the findings of this review are available as follows:
(i)~the \textbf{complete search strings} per database are registered in
supplementary file~S1 and in the project's public repository
([URL removed for blind review]);
(ii)~the \textbf{standardised data extraction form} (fields extracted per
study, including authors, year, study design, population, technology used,
outcomes and methodological quality assessment) is available as supplementary
material~S2;
(iii)~\textbf{screening decisions} (inclusion/exclusion by title, abstract
and full text), with the respective reasons for exclusion, were recorded in
the Rayyan platform \citep{Ouzzani2016} and are made available upon request
to the corresponding author;
(iv)~the \textbf{complete list of included studies}, with classification by
thematic axis and quality score, is incorporated in Table~\ref{tab:estudos}
in the body of the article.
Raw data from recordings, transcriptions or patient records do not apply to
this study, which is based exclusively on published literature.

\subsection*{Use of artificial intelligence}
In accordance with Elsevier's Generative AI Policy, the authors declare that
artificial intelligence tools were used in the preparation of this work for
literature search assistance, corpus screening support, data extraction and
translation of excerpts from articles consulted in a foreign language.
All intellectual content, analyses, syntheses and conclusions are the sole
responsibility of the authors, who fully reviewed and validated all material
produced with AI assistance.

%% ================================================================
\bibliographystyle{elsarticle-num}
\bibliography{references}

\end{document}
